\documentclass[11pt]{article}
\usepackage[utf8]{inputenc}
\usepackage[margin=1in]{geometry}
\usepackage{amsmath}
\usepackage{graphicx}
\usepackage{booktabs}
\usepackage{hyperref}
\usepackage{natbib}
\usepackage{setspace}
\onehalfspacing

\title{Shape-Changing Electrode Array for Minimally Invasive Large-Scale Intracranial Brain Activity Mapping}
\author{Author A$^{1}$, Author B$^{1,2}$, Author C$^{2}$ \\
\small $^{1}$Department of Biomedical Engineering, University of Research \\
\small $^{2}$Department of Neurosurgery, Medical Center}
\date{}

\begin{document}
\maketitle

\begin{abstract}
Large-scale brain activity mapping with high spatiotemporal resolution requires electrocorticographic (ECoG) electrode arrays, but conventional implantation demands extensive craniotomy and durotomy with significant surgical risks. Here we present a shape-changing electrode array (SCEA) that transforms from a narrow compressed strip (3 mm wide) into a large conformal sheet (up to 20 mm $\times$ 15 mm) upon reaching body temperature, enabling implantation through a small skull opening. The SCEA uses a 4 $\mu$m thick Parylene-C substrate with carbon nanotube (CNT)/gold multilayer electrodes that maintain conductivity under extreme mechanical deformation, coupled with a nitinol shape-memory alloy actuator that transitions from martensite to austenite at 37$^{\circ}$C. We demonstrate stable impedance ($\sim$50 k$\Omega$ at 1 kHz) before and after shape transformation, no brain structural changes on MRI up to 8 weeks post-implantation in rats, and minimal inflammatory response on immunohistochemistry comparable to standard epidural implants. The 64-channel SCEA recorded high-quality micro-ECoG during 4-aminopyridine-induced seizures in rats, capturing high-frequency oscillations up to 500 Hz and spatiotemporal propagation patterns. A 256-channel SCEA was implanted subdurally in canine (beagle) brains through a small dura slit, successfully mapping cortical activity during emergence from anesthesia. Post-operative MRI-based localization enabled precise electrode position mapping. The SCEA provides a minimally invasive approach to large-scale intracranial brain activity mapping.
\end{abstract}

\section{Introduction}

Understanding brain function requires recording neural activity at high spatial and temporal resolution across large cortical areas. Among available techniques, electrocorticography (ECoG) provides a unique combination of millimeter-scale spatial resolution, millisecond temporal resolution, and access to high-frequency activity up to 500 Hz, including high-gamma oscillations and ripple frequencies relevant to cognition and pathological states \citep{muller2016,engel2005}. However, conventional ECoG array implantation requires an extensive craniotomy (removal of a large bone flap) and durotomy (opening of the dura mater), procedures that carry significant risks including infection, inflammation, cerebral edema, and hemorrhage \citep{vanmiert2020,wong2009}.

The surgical invasiveness of ECoG limits its clinical application primarily to presurgical epilepsy monitoring, where the diagnostic benefit outweighs the procedural risks. For research applications and emerging brain-computer interfaces, the invasiveness of ECoG has driven interest in less invasive alternatives, including scalp EEG, magnetoencephalography (MEG), and hemodynamic imaging (fMRI, fNIR). However, each of these modalities involves fundamental trade-offs in spatial resolution, temporal resolution, or bandwidth (Table~\ref{tab:comparison}).

\begin{table}[h]
\centering
\caption{Comparison of brain mapping techniques.}
\label{tab:comparison}
\begin{tabular}{lcccc}
\toprule
Technique & Spatial Res. & Temporal Res. & Bandwidth & Invasiveness \\
\midrule
Scalp EEG & cm & ms & 0.5--50 Hz & Non-invasive \\
MEG & mm & ms & Up to 250 Hz & Non-invasive \\
fMRI & mm & s & Hemodynamic & Non-invasive \\
fNIR & mm--cm & s & Hemodynamic & Non-invasive \\
Standard ECoG & $\mu$m--mm & ms & Up to 500 Hz & Highly invasive \\
SCEA (this work) & $\mu$m--mm & ms & Up to 500 Hz & Minimally invasive \\
\bottomrule
\end{tabular}
\end{table}

An ideal solution would achieve the spatial resolution, temporal resolution, and bandwidth of ECoG through a minimally invasive implantation procedure. This requires an electrode array that can be inserted through a small opening but then expand to cover a large cortical area. Several approaches have been explored, including inflatable subdural electrodes and robotic insertion systems, but none has achieved the combination of large area coverage, high channel count, and minimal surgical footprint needed for clinical translation \citep{rogers2010,viventi2011}.

Here we present the shape-changing electrode array (SCEA), which exploits the shape-memory effect of nitinol (nickel-titanium alloy) to transform from a compressed strip that fits through a 3 mm slit to a full-size conformal sheet covering over 300 mm$^{2}$ of cortical surface. The SCEA uses ultrathin (4 $\mu$m) Parylene-C substrates with CNT/Au multilayer electrodes that maintain electrical function throughout the extreme mechanical deformation of the shape transformation.

\section{Methods}

\subsection{SCEA Design and Fabrication}

The SCEA comprises multiple functional layers (Table~\ref{tab:specs}). The substrate is a 4 $\mu$m thick Parylene-C film deposited by chemical vapor deposition, providing excellent flexibility and biocompatibility. SU8 photoresist insulation layers define the electrode openings and interconnect traces. The recording sites (100 $\times$ 100 $\mu$m$^{2}$ each) consist of a CNT/Au bilayer: a web-like single-walled CNT thin film grown by chemical vapor deposition, overlaid with a thin gold layer for optimized impedance.

\begin{table}[h]
\centering
\caption{SCEA design specifications.}
\label{tab:specs}
\begin{tabular}{ll}
\toprule
Parameter & Value \\
\midrule
Substrate material & Parylene-C \\
Substrate thickness & 4 $\mu$m \\
Insulation layers & SU8 \\
Conductive layer & CNT/Au (SWCNT by CVD) \\
Active site size & 100 $\times$ 100 $\mu$m$^{2}$ \\
Compressed strip width & 3 mm \\
Deployed array size (rat) & 20 mm $\times$ 15 mm \\
Channel count (rat) & 64 \\
Deployed array size (canine) & 20 mm $\times$ 10 mm \\
Channel count (canine) & 256 \\
Shape actuator & Nitinol (NiTi) \\
Actuator transition temperature & 37$^{\circ}$C \\
Bonding adhesive & PEO (water-soluble) \\
\bottomrule
\end{tabular}
\end{table}

The CNT layer is critical for maintaining electrical function during shape transformation. While nanometer-thick gold films crack and lose conductivity under large strain, the web-like structure of the CNT network accommodates extreme deformation by rearranging nanotube contacts without breaking individual tubes. The CNT/Au hybrid thus preserves conductivity throughout the compression and expansion cycle.

\subsection{Shape-Memory Actuator}

The nitinol actuator is pre-shaped at high temperature into a concave polygon matching the desired deployed array geometry, then cooled below the martensite transition temperature and mechanically compressed into a narrow strip. At body temperature (37$^{\circ}$C), the nitinol transitions to its austenite phase and recovers the pre-shaped geometry, deploying the electrode array. The actuator is temporarily bonded to the array using polyethylene oxide (PEO) adhesive, which dissolves in aqueous solution after deployment, allowing the actuator to be withdrawn while the electrode array remains in place.

\subsection{Implantation Procedure}

For rat experiments (epidural), a small slit was made in the skull. The compressed SCEA strip was inserted through the slit onto the brain surface. Upon warming to body temperature, the nitinol actuator deployed the array into its full sheet configuration. The actuator was then withdrawn, and the electrode array remained conformally on the cortical surface.

For canine experiments (subdural), a small slit was made in the dura after a minimal craniotomy. The 256-channel SCEA was inserted subdurally and deployed using the same shape-memory mechanism. Deployment was monitored visually to ensure no vascular disruption.

\subsection{In Vitro Characterization}

Electrochemical impedance spectroscopy was performed before and after one compression/expansion cycle in phosphate-buffered saline. Impedance was measured at 1 kHz for all channels.

\subsection{MRI Studies in Rats}

Longitudinal T2-weighted structural MRI was performed at 1, 2, 3, 4, and 8 weeks post-implantation to monitor for brain structural changes. Dynamic contrast-enhanced MRI (DCE-MRI) was used to assess blood-brain barrier (BBB) integrity and meningeal enhancement. A control group received the same electrode through a standard craniotomy (without shape transformation) for comparison.

\subsection{Immunohistochemistry in Rats}

Rats were perfused at 1, 4, and 8 weeks post-implantation. Coronal sections through the implant area were stained for GFAP (astrocyte reactivity), Iba1 (microglial activation), and NeuN (neuronal density). Cortex was divided into L1, upper layers (UL), and lower layers (LL) for quantification. The contralateral hemisphere served as the within-animal control.

\subsection{Rat Seizure Recording}

Epileptiform activity was induced by topical application of 4-aminopyridine (4-AP) through a small cranial hole. Epidural micro-ECoG was recorded with the 64-channel SCEA at broadband bandwidth (0.5--500 Hz).

\subsection{Canine Emergence Recording}

Following subdural implantation in beagle dogs, ECoG was recorded during emergence from general anesthesia. The 256-channel array provided spatial mapping of cortical activity changes across the transition from deep anesthesia through light anesthesia to wakefulness.

\subsection{Post-Operative Localization}

For canine experiments, post-operative T1-weighted MRI with nickel-enhanced SCEA (Ni-SCEA) markers enabled visualization of electrode positions. Three-dimensional reconstruction using MIPAV software mapped individual channel positions onto the cortical surface.

\section{Results}

\subsection{Impedance Stability Through Shape Transformation}

The 64-channel SCEA maintained stable impedance through the compression and expansion cycle (Table~\ref{tab:impedance}). Mean impedance at 1 kHz was approximately 50 k$\Omega$ both before and after transformation, with low channel-to-channel variation. All 64 channels remained functional after transformation, demonstrating that the CNT/Au electrode structure withstands the mechanical deformation without degradation.

\begin{table}[h]
\centering
\caption{Impedance characterization of 64-channel SCEA.}
\label{tab:impedance}
\begin{tabular}{lcc}
\toprule
Measurement & Before Transform & After Transform \\
\midrule
Mean impedance (1 kHz) & $\sim$50 k$\Omega$ & $\sim$50 k$\Omega$ \\
Channel-to-channel variation & Low & Low \\
Functional channels & 64/64 & 64/64 \\
\bottomrule
\end{tabular}
\end{table}

\subsection{MRI Shows No Brain Structural Changes}

Longitudinal T2-weighted MRI from 1 to 8 weeks post-implantation revealed no structural changes in the brain parenchyma beneath the SCEA. DCE-MRI showed segments of meningeal convexity enhancement at 1 week post-implantation, which progressively diminished and completely resolved by 4 weeks. Critically, no cortical BBB breach was observed at any timepoint. The meningeal enhancement level was comparable between the minimally invasive SCEA implantation and the standard craniotomy control, indicating that the shape transformation itself does not introduce additional tissue disruption.

\subsection{Minimal Inflammatory Response on IHC}

Immunohistochemical analysis revealed mild increases in GFAP and Iba1 staining at 1 week post-implantation, consistent with the expected acute surgical response. By 4 weeks, both markers had returned to near-baseline levels, and by 8 weeks, the tissue beneath the SCEA was indistinguishable from the contralateral control. NeuN staining showed no changes in neuronal density at any timepoint, confirming that the SCEA does not cause neuronal loss. The inflammatory response was comparable to that seen with standard epidural electrode implants, demonstrating that the minimally invasive deployment does not compromise biocompatibility.

\subsection{High-Quality Seizure Recording in Rats}

The 64-channel SCEA recorded high-quality micro-ECoG signals during 4-AP-induced seizures in rats. Low-frequency oscillations (0.5--50 Hz), ripples (80--250 Hz), and fast ripples (250--500 Hz) were all clearly resolved. Spatiotemporal propagation of seizure activity was mapped across the array, demonstrating the spatial coverage and temporal resolution of the SCEA. The seizure wavefront could be tracked across individual channels, providing information about propagation speed and direction that would not be available from lower-density recordings.

\subsection{Subdural Canine Implantation and Recording}

The 256-channel SCEA was successfully implanted subdurally in beagle dogs through a small dura slit. The shape transformation occurred smoothly upon warming to body temperature, with the array deploying to its full 20 mm $\times$ 10 mm extent without disrupting cortical vasculature. The deployed array conformed to the brain surface curvature.

During emergence from anesthesia, the SCEA captured the progressive transition from suppressed, globally synchronized activity (burst-suppression pattern under deep anesthesia) through increasing regional differentiation to active, spatiotemporally organized cortical patterns characteristic of the awake state. The full spatial coverage of the 256-channel array enabled simultaneous monitoring of activity across a large cortical territory, revealing that emergence-related changes were not spatially uniform but showed region-specific temporal dynamics.

\subsection{Post-Operative Electrode Localization}

MRI-based localization of the Ni-SCEA in the canine brain enabled precise mapping of each electrode position to the underlying cortical anatomy. Serial MRI sections through the implant area were reconstructed in three dimensions using MIPAV software, providing a complete picture of the array position relative to cortical sulci and gyri. This localization capability is essential for relating recorded signals to specific cortical regions.

\section{Discussion}

The SCEA represents a new paradigm for intracranial brain activity mapping: achieving the spatial resolution, temporal resolution, and bandwidth of conventional ECoG while dramatically reducing surgical invasiveness. The key innovation is the use of a shape-memory actuator that transforms the array from a narrow strip that fits through a minimal skull opening to a large conformal sheet covering hundreds of square millimeters of cortical surface.

The CNT/Au multilayer electrode structure is critical to this approach. Conventional thin-film metal electrodes crack under the strains imposed by the shape transformation, but the web-like CNT network accommodates deformation by redistributing inter-tube contacts without breaking individual nanotubes \citep{lipomi2011,khang2006}. The stability of impedance through the transformation cycle confirms that the CNT layer preserves electrical function.

The biocompatibility data are encouraging for clinical translation. The absence of brain structural changes on MRI, the resolution of meningeal enhancement by 4 weeks, and the minimal inflammatory response on IHC all suggest that the SCEA is well tolerated over chronic implantation periods. The comparable inflammatory response between minimally invasive and standard craniotomy implantation confirms that the shape transformation does not introduce additional tissue damage.

The demonstration of subdural implantation in a large animal model (beagle) represents an important step toward clinical translation. The canine brain is sufficiently large to require the SCEA's full spatial coverage, and the successful recording during anesthesia emergence demonstrates that the device can capture clinically relevant brain state transitions. The 256-channel design provides the channel density needed for detailed spatial mapping of cortical activity.

Several challenges remain for clinical translation. The current design requires a small craniotomy and dura slit, which, while much less invasive than conventional approaches, still requires a surgical procedure under general anesthesia. The nitinol actuator must be carefully designed to match the target cortical anatomy, and the deployment process must be monitored to avoid vascular damage. Long-term chronic recording stability beyond 8 weeks needs to be established. Finally, clinical-grade fabrication processes must be developed to ensure reproducibility and reliability \citep{rogers2010}.

\section{Conclusion}

We have developed a shape-changing electrode array that enables minimally invasive large-scale intracranial brain activity mapping. The SCEA transforms from a 3 mm compressed strip to a large conformal sheet at body temperature, maintaining stable electrode impedance throughout the transformation. Comprehensive biocompatibility testing in rats demonstrates no brain structural changes or significant inflammatory response. High-quality broadband micro-ECoG recording was demonstrated in both rat seizure models and canine emergence paradigms. The SCEA provides a practical approach to achieving the high spatiotemporal resolution of ECoG while substantially reducing surgical invasiveness.

\bibliographystyle{plainnat}
\bibliography{references}

\end{document}
